
\section{Materials and Methods}
\label{sec:materials_and_methods}
This study builds upon data previously collected in the work of Sanghan et al.~\cite{ref22}, which involved a comprehensive investigation comparing the 3D 
shank kinematics of stroke patients to healthy individuals.

\subsection{Data Source}
The original experiment included sixteen stroke patients and sixteen healthy individuals, all of whom walked on a flat 10-meter indoor walkway while wearing IMU sensors 
(MBIENTLAB Inc., USA) placed bilaterally on their shanks as seen on Figure \ref{fig:sensor_placement}. Data acquisition protocols, participant criteria, and ethical 
approvals were thoroughly detailed in the original publication~\cite{ref22}. The recorded data consisted of tri-axial gyroscope signals sampled at 100 Hz, with angular 
velocities measured along the sagittal (x), frontal (y), and transverse (z) planes.


\subsection{Feature Engineering}
This study presents a structured pipeline for extracting gait cycle features from gyroscopic data collected using Inertial Measurement Units (IMUs) placed 
on the left and right shanks. The analysis focuses exclusively on the z-axis angular velocity component, as it effectively captures sagittal plane rotational 
movement, making it a strong indicator of lower-limb kinematics during gait. \\

\textbf{Gait Cycle Detection}
Let $G_L(t) \in \mathbb{R}$ denote the filtered z-axis gyroscope signal from the left shank over time $t$. To identify gait cycles, we apply peak detection 
to $G_L(t)$ based on amplitude and temporal separation constraints. Peaks are defined as local maxima $P = \{p_1, p_2, \dots, p_n\}$ such that: 

\[G_L(p_i) > \theta \quad \text{and} \quad p_{i+1} - p_i \geq d,\]

where $\theta = 0.5$ deg/s and $d = 80$ samples (at 100 Hz) to ensure biomechanical plausibility. Each gait cycle is defined by the interval $[p_i, p_{i+1})$. 
Only intervals with a minimum width of 10 samples are retained. \\

\textbf{Feature Computation}
Within each valid gait cycle window, we compute a comprehensive set of features from both the left $G_L(t)$ and right $G_R(t)$ shank signals.
\begin{enumerate}
    \item Descriptive Statistics: including mean, standard deviation, maximum and minimum values of the gyroscope signal. \\ 
    \item Motion Score: We define a custom metric reflecting signal amplitude dynamics:
    \[
    M = \max \left( |G(t)| \right) - \min \left( |G(t)| \right).
    \]
    \newpage
    \item Stance and Swing Phase Detection: Phase segmentation is performed using zero-crossing detection on the filtered gyroscope signal. Given zero-crossings at 
    time indices $z_i$, we define valid intervals $[z_i, z_{i+1}]$. The minimum angular velocity within each interval marks the transition between stance and swing. 
    The duration of each phase is computed as: 
    \[
    T_{\text{stance}} = t_{\min} - t_{\text{start}},
    \]
    \[
    T_{\text{swing}} = t_{\text{end}} - t_{\min},
    \]
    where $t_{\min}$ is the timestamp at which the minimum occurs between zero-crossings. \\

    \item Stride Time: The stride time for each leg is computed as:
    \[
    T_{\text{stride}} = T_{\text{stance}} + T_{\text{swing}}
    \]
\end{enumerate}


\textbf{Asymmetry Metrics and Label Assignment}
To quantify gait asymmetry, we compute two indices based on the stride durations of the left ($T_L$) and right ($T_R$) limbs:

\begin{enumerate}
    \item Asymmetry Index (AI):
    \[
    \text{AI} = \frac{T_L - T_R}{T_L + T_R}, \quad T_L + T_R \neq 0.
    \]
    
    \item Symmetry Ratio (SR):
    \[
    \text{SR} = \frac{\min(T_L, T_R)}{\max(T_L, T_R)}, \quad \max(T_L, T_R) \neq 0 \\
    \]
\end{enumerate}

If either $T_L$ or $T_R$ is undefined for a given cycle, the window is marked invalid, and no asymmetry metrics are computed. Following this rule we were able to
obtain a total of over 2500 valid gait cycles from the original dataset as seen on the Table~\ref{tab:valid_invalid_cycles}. 
\begin{table}[h]
    \centering
    \caption{Number of valid and invalid gait cycles}
    \label{tab:valid_invalid_cycles}
    \renewcommand{\arraystretch}{1.2}
    \begin{tabular}{cc}
    \toprule
    \textbf{Valid Gait Cycles} &\textbf{Invalid Gait Cycles} \\
    \toprule
        2837 & 884 \\
    \bottomrule
    \end{tabular}
\end{table}

\newpage
Using the two symmetry metrics, Asymmetry Index and Symmetry Ratio, each gait window is categorized into a confidence level that reflects the likelihood and 
severity of gait asymmetry. These thresholds are designed to support clinical interpretability and model stratification based on asymmetry certainty.

\textbf{Confidence Levels}
Table~\ref{tab:confidence_distribution} summarizes the distribution of gait windows categorized into the three confidence levels. As expected, a greater proportion 
of the high-confidence asymmetric cycles belong to the stroke group, whereas the healthy class dominates the low-confidence windows. The moderate-confidence 
group reflects a transitional state and includes a mixed distribution of both subject types. Despite the overall class imbalance, especially at high confidence 
levels, stroke-related asymmetry remains consistently detectable across all tiers. This stratification confirms that even patients undergoing rehabilitation may 
occasionally exhibit gait cycles resembling normal patterns.


\begin{table}[h]
    \centering
    \caption{Distribution of Stroke vs. Healthy Gait Windows by Confidence Level}
    \label{tab:confidence_distribution}
    \renewcommand{\arraystretch}{1.2}
    \begin{tabular}{lccc}
    \toprule
    \textbf{Confidence Level} & \textbf{Class Rules} & \textbf{Stroke Windows} & \textbf{Healthy Windows} \\
    \midrule
    High      & $\text{AI} > 0.2$ or $\text{SR} < 0.8$   & 586  & 2251 \\
    Moderate  & $\text{AI} > 0.15$ or $\text{SR} < 0.75$ & 484  & 2353  \\
    Low       & $\text{AI} > 0.1$ or $\text{SR} < 0.9$   & 367  & 2470   \\
    \bottomrule
    \end{tabular}
\end{table}


\subsection{Feature Selection}
We conducted Recursive Feature Elimination (RFE) analysis using a Random Forest Classifier to refine the feature set further. In this process we 
iteratively removed the least important features based on the contribution to the accuracy of the model, until we got the optimal subset of features.
The RFE results, shown in Table~\ref{tab:rfe}, indicate that the most significant features for classification include stance and swing times for both legs,
as well as the motion scores derived from the gyroscope signals. The stance and stride times were consistently ranked as the most important across all confidence 
levels, while the left motion score was selected for the high and moderate confidence levels, and the right motion score for the low confidence level.

\begin{table}[h]
    \centering
    \caption{Recursive feature elimination (RFE) results}
    \small
    \begin{tabular}{l c c}
        \hline
        \textbf{Feature} & \textbf{Importance} & \textbf{Ranking} \\
        \hline
        \texttt{right\_stance\_time}  & 0.18 & 1 \\
        \texttt{right\_swing\_time}   & 0.12 & 1 \\
        \texttt{left\_stance\_time}   & 0.11 & 1 \\
        \texttt{left\_swing\_time}    & 0.05 & 1 \\
        \texttt{left\_motion\_score}  & 0.05 & 1 \\
        \texttt{right\_motion\_score} & 0.05 & 1 \\
        \hline
    \end{tabular}
    \label{tab:rfe}
\end{table}

Provided that we wanted to keep the model interpretable and medically relevant, which means that we are going to use the stride and stance times as
the main features for the classification task, since those features can be easily computed in real-time using IMUs and smartphone solution.

\subsection{Classification Models}
The relatively small dataset size along with the modest dimensionality of the feature space, guided us to use classical machine learning models rather than 
deep learning approaches. Deep learning models typically require large volumes of training data to generalize effectively and are more prone to overfitting 
in small-data problems. On the other hand, machine learning models are well-suited for structured datasets with limited size, providing robust performance 
and greater interpretability. To evaluate classification performance, we tested a diverse set of algorithms: Logistic Regression, Random Forest, Support 
Vector Machines (SVM) with linear, radial basis function (RBF), and polynomial kernels, Extreme Gradient Boosting (XGBoost), and K-Nearest Neighbors (KNN). 
These models offer a mix of linear and non-linear decision boundaries, ensemble learning, and instance-based reasoning, enabling a comprehensive comparison 
across different learning paradigms.
